\documentclass{article}
\usepackage{spconf,amsmath,graphicx,booktabs,url,amssymb}

\ninept

\title{Retrieval in Pair-Difference Space and Transductive System Smoothing \\
       for Speaker and Accent Similarity Prediction}

\name{Ranjit Patro}
\address{Independent Researcher, India \\
         \texttt{ranjitpatro100@gmail.com}}

\begin{document}
\maketitle

\begin{abstract}
We describe our submission to Track 3 of the VoiceMOS Challenge 2026, which asks
for utterance-level prediction of speaker similarity (\textsc{spk}) and accent
similarity (\textsc{acc}) between a codec- or TTS-generated utterance and a
natural reference recording of the target speaker reading a \emph{different}
sentence. Our system uses no backbone fine-tuning: it fuses cosine features from
seven frozen pretrained models with a retrieval component, a small trained
pairwise head, and a system-level prior. Its two distinguishing components are
(i) \emph{retrieval in pair-difference space}, a variant of RAMP in which the
retrieval key is the difference between the two utterances' self-supervised
embeddings rather than a single utterance, and (ii) \emph{transductive
system-mean smoothing}, a label-free inference-time step that pulls each
prediction toward the predicted mean of its own system. On the official
evaluation set our system reaches 0.549 / 0.474 utterance-level SRCC and
0.942 / 0.902 system-level SRCC for \textsc{spk} / \textsc{acc}, placing third
of seven teams at system level and improving on both official baselines by
0.14 and 0.09 SRCC respectively. We further report several negative results,
including a reproducible failure of listener-dependent modelling on this data.
\end{abstract}

\keywords{speech quality assessment, speaker similarity, accent similarity,
retrieval augmentation, self-supervised speech models}

\section{Introduction}

Subjective evaluation remains the reference standard for assessing generated
speech, and automatic mean-opinion-score (MOS) prediction aims to approximate it
cheaply \cite{utmos,mosbench}. Track 3 of the VoiceMOS Challenge 2026
\cite{codecmos} departs from single-utterance MOS in two respects: the quantity
being predicted is a \emph{similarity} between two recordings, and it is split
into two correlated sub-scores --- speaker similarity and accent similarity.
Because the reference utterance is a different sentence spoken by the target
speaker, the task cannot be reduced to signal-level comparison and must rely on
content-independent speaker and accent representations.

Two properties of the data shape any reasonable solution. First, the training
set is small (2{,}800 rated pairs) and human ratings are noisy, which strongly
favours frozen pretrained representations over fine-tuning; this echoes findings
from earlier challenge editions where compact models outperformed large
fine-tuned ones \cite{vmc23}. Second, both the development and evaluation sets
contain systems that never appear in training --- 57\% of evaluation pairs come
from four unseen systems --- so a model that implicitly memorises system
identity will not generalise.

Our contributions are: (1) a retrieval-augmented predictor that retrieves
neighbouring \emph{pairs} in the space of embedding differences, with an
ablation isolating why the retrieval space matters; (2) a system-level prior
with an out-of-distribution-safe fallback, together with a leave-one-system-out
stress test showing it never degrades performance; (3) transductive system-mean
smoothing, a label-free inference-time step that provided the single largest
gain on the evaluation set; (4) a grouped cross-validation harness calibrated
against the public leaderboard; and (5) a set of negative results, including
evidence that listener-dependent training hurts on this dataset.

\section{Task and Data}

Each item is a pair $(a,b)$ where $a$ is the output of a codec-resynthesis or
neural-codec TTS system and $b$ is a natural VCTK recording of the target
speaker reading a different sentence (3--7\,s, 16\,kHz). The targets are two
utterance-level MOS values in $[1,5]$, and the primary metric is the
utterance-level Spearman rank correlation coefficient (SRCC).

The corpus \cite{codecmos} contains 4{,}000 samples from 24 systems covering 32
VCTK speakers and 10 English accents. The training partition has 2{,}800 pairs
from 21 systems with 13{,}687 listener-wise ratings (mean 4.9 raters per pair);
the development and evaluation partitions each contain 600 pairs, with 2 and 4
unseen systems respectively.

Three statistics are relevant to the modelling choices below. The two sub-scores
correlate at $0.86$ at utterance level and $0.79$ even within a single system,
so they are not independently identifiable from these data. Within-pair listener
standard deviation is $0.85$ (\textsc{spk}) and $0.92$ (\textsc{acc}) on a
five-point scale; a split-half reliability estimate places the attainable
utterance-level correlation near $0.81$ and $0.76$ respectively, so utterance
level SRCC is intrinsically capped well below unity. Finally, between-system
variance accounts for 45\% (\textsc{spk}) and 33\% (\textsc{acc}) of total
variance, which motivates exploiting system-level structure explicitly.

\section{Proposed System}

\subsection{Frozen feature panel}

We extract, for every waveform, representations from seven public pretrained
models used entirely frozen: ECAPA-TDNN (192-d), TitaNet-Large (192-d),
WeSpeaker ResNet34 (256-d), CommonAccent-ECAPA \cite{commonaccent} (192-d
embedding plus a 16-class posterior), a wav2vec2 accent classifier (768-d hidden
mean plus a 13-class posterior), WavLM-Large \cite{wavlm} (all 25 hidden layers,
mean-pooled over time), and the UTMOS \cite{utmos} scalar quality score.
Waveforms are used at 16\,kHz mono with each model's own front end; no
augmentation, trimming or volume normalisation is applied.

From these we build roughly 35 scalar pair features: the cosine similarity
between $a$ and $b$ for every embedding source, the per-layer cosine for all 25
WavLM layers, and $\mathrm{UTMOS}(a)$, $\mathrm{UTMOS}(b)$ and their difference.
A ridge regressor ($\alpha=1$) on standardised features gives the parametric
prediction $\hat{y}_{\text{ridge}}$. Scalar features cannot encode system
identity, which is deliberate.

\subsection{Retrieval in pair-difference space}

RAMP \cite{ramp} augments a parametric MOS predictor with a $k$-nearest-neighbour
retrieval over a training datastore, which is attractive in low-resource,
out-of-distribution settings. The natural adaptation to a \emph{pairwise} task is
to make the retrieval key describe the pair, not one utterance. We therefore
define the pair-difference vector
\begin{equation}
  \mathbf{d}(a,b) = \bigoplus_{m \in \mathcal{M}}
    \left( \frac{\mathbf{e}_m(a)}{\lVert \mathbf{e}_m(a)\rVert}
         - \frac{\mathbf{e}_m(b)}{\lVert \mathbf{e}_m(b)\rVert} \right),
\end{equation}
where $\bigoplus$ is concatenation and $\mathcal{M}$ comprises ECAPA,
CommonAccent, WavLM layer 3 (speaker-dominant) and WavLM layer 12
(accent-dominant), giving 2{,}432 dimensions. We standardise, reduce to 128
dimensions with PCA, and retrieve the $k=30$ nearest training pairs, weighting
neighbour labels by $\mathrm{softmax}(-d_i/\bar{d})$ to obtain
$\hat{y}_{\text{ret}}$. The two predictions are fused in $z$-space,
\begin{equation}
  \hat{y}_{\text{fuse}} = \beta\, z(\hat{y}_{\text{ridge}})
                        + (1-\beta)\, z(\hat{y}_{\text{ret}}),
\end{equation}
with $\beta=0.8$ (\textsc{spk}) and $0.7$ (\textsc{acc}).

\subsection{Trained pairwise head}

For ensemble diversity we train a small head on frozen WavLM features. Per-task
learnable softmax weights collapse the 25 layers into a single 1024-d vector,
which a linear layer projects to 128-d. The pair representation
$[\mathbf{p}_a, \mathbf{p}_b, |\mathbf{p}_a - \mathbf{p}_b|,
\mathbf{p}_a \odot \mathbf{p}_b]$ feeds a shared trunk and two output heads. The
loss combines mean squared error with a margin-ranking term (margin $0.3$)
computed over batches sampled \emph{across} systems, so that the ranking
objective cannot be satisfied by system identity alone. The head has 164{,}212
trainable parameters; we average six seeds and combine with
$\hat{y}_{\text{fuse}}$ by equally weighted $z$-averaging.

\subsection{System-level prior with OOD-safe fallback}

Because between-system variance dominates, we blend the prediction toward the
training-label mean of the pair's system,
$\hat{y}_{\text{prior}} = \alpha\, z(\hat{y}) + (1-\alpha)\, z(\bar{y}_{s})$,
with $\alpha=0.5$ (\textsc{spk}) and $0.6$ (\textsc{acc}). For a system unseen in
training no prior exists, and the model falls back to the plain cosine-based
prediction. This fallback is what makes the prior safe under distribution shift
(Section~\ref{sec:ablation}).

\subsection{Transductive system-mean smoothing}

The evaluation manifest reveals which system each pair belongs to, though not its
quality. Since system-level signal is far more reliable than utterance-level
signal, we apply a label-free smoothing step at inference,
\begin{equation}
  \hat{y}' = w\,\hat{y} + (1-w)\,\frac{1}{|S|}\sum_{j \in S} \hat{y}_j ,
  \label{eq:smooth}
\end{equation}
where $S$ is the set of evaluation pairs from the same system and $w=0.7$. This
uses no labels and is applied identically to seen and unseen systems. It
sharpens the between-system ordering, which dominates SRCC, at the cost of
within-system resolution that is largely rater noise. Scores are finally rescaled
to $[1,5]$.

\section{Validation Methodology}

Because unseen systems dominate the evaluation set, we select all
hyper-parameters using \emph{grouped cross-validation by system}: seven folds,
each holding out whole systems. This choice was not cosmetic. A naive random
split reported $0.63$ SRCC for an early model whose true development score was
$0.36$, because random splits let the model memorise system identity.

We verified that the harness tracks the leaderboard before trusting it. The
zero-shot ECAPA cosine baseline scores $0.432$ under grouped CV, exactly matching
its official development score of $0.432$; a wav2vec2 posterior cosine scores
$0.323$ against $0.319$ on the leaderboard. This calibration allowed design
decisions to be made locally, without consuming submissions.

\section{Results}
\label{sec:results}

\subsection{Official evaluation results}

Table~\ref{tab:eval} reports our official Track 3 results together with the two
organiser baselines. Our system improves on the stronger baseline by $0.135$
(\textsc{spk}) and $0.088$ (\textsc{acc}) utterance-level SRCC. At system level
it reaches $0.942$ / $0.902$, placing third of seven teams for both sub-scores
and within $0.006$ of the best system-level speaker score in the challenge.

\begin{table}[t]
\centering
\caption{Official evaluation results (SRCC). Ranks are among the seven
participating teams; B1/B2 are the organiser baselines.}
\label{tab:eval}
\small
\begin{tabular}{lcccc}
\toprule
 & \multicolumn{2}{c}{Utterance-level} & \multicolumn{2}{c}{System-level} \\
\cmidrule(lr){2-3}\cmidrule(lr){4-5}
System & \textsc{spk} & \textsc{acc} & \textsc{spk} & \textsc{acc} \\
\midrule
B1 (ECAPA cosine)      & 0.414 & 0.386 & 0.874 & 0.838 \\
B2 (ECAPA + head)      & 0.403 & 0.316 & 0.914 & 0.786 \\
Best team              & 0.644 & 0.565 & 0.948 & 0.943 \\
\midrule
\textbf{Ours}          & \textbf{0.549} & \textbf{0.474}
                       & \textbf{0.942} & \textbf{0.902} \\
\quad rank             & 4th & 5th & 3rd & 3rd \\
\bottomrule
\end{tabular}
\end{table}

\subsection{Effect of transductive smoothing}

Smoothing (Eq.~\ref{eq:smooth}) was validated on 3{,}400 labelled pairs under
grouped CV, where it gained $+0.008$ for both sub-scores, with a broad optimum
over $w\in[0.5,0.9]$ rather than a sharp peak. On the withheld evaluation set the
gain was substantially larger, as shown in Table~\ref{tab:smooth}: $+0.037$
(\textsc{spk}) and $+0.016$ (\textsc{acc}). We attribute the amplification to the
composition of the evaluation set, in which two unseen systems contribute 160
pairs each; such large groups yield very stable predicted means, exactly the
regime in which Eq.~\ref{eq:smooth} is most effective.

\begin{table}[t]
\centering
\caption{Transductive system-mean smoothing, evaluation set (SRCC).}
\label{tab:smooth}
\small
\begin{tabular}{lccc}
\toprule
System & \textsc{spk} & \textsc{acc} & mean \\
\midrule
Base (retrieval + head + prior) & 0.512 & 0.458 & 0.485 \\
\quad + smoothing ($w=0.7$)     & \textbf{0.549} & \textbf{0.474} & \textbf{0.512} \\
$\Delta$                        & $+0.037$ & $+0.016$ & $+0.026$ \\
\bottomrule
\end{tabular}
\end{table}

\subsection{Ablations}
\label{sec:ablation}

Table~\ref{tab:ablate} isolates the retrieval space under grouped CV. Retrieval
over the compact cosine scalars is statistically indistinguishable from ridge
alone; only retrieval in the high-dimensional embedding-difference space helps.
A single shared retrieval space outperforms per-target spaces, consistent with
the $0.86$ coupling between the sub-scores, and four well-chosen sources beat
nine, indicating that redundant embeddings dilute the neighbour distance.

\begin{table}[t]
\centering
\caption{Retrieval-space ablation (grouped CV, pooled out-of-fold SRCC).}
\label{tab:ablate}
\small
\begin{tabular}{lcc}
\toprule
Retrieval space & \textsc{spk} & \textsc{acc} \\
\midrule
Ridge only (no retrieval)            & 0.644 & 0.598 \\
Cosine scalars (35-d)                & 0.648 & 0.602 \\
\textbf{Embedding difference, 4 sources} & \textbf{0.665} & \textbf{0.630} \\
Embedding difference, per-target     & 0.653 & 0.623 \\
Embedding difference, 9 sources      & 0.659 & 0.621 \\
\bottomrule
\end{tabular}
\end{table}

We also stress-tested the system prior by sweeping the fraction of unseen systems
in a leave-one-system-out protocol. The prior's advantage over the plain cosine
prediction decreases monotonically from $+0.010$ at 50\% unseen to $+0.001$ at
90\% unseen, but never becomes negative. The fallback design therefore makes the
prior safe under severe distribution shift, in contrast to reports that learned
domain priors degrade out of domain.

\section{Negative Results}

We report four negative findings, each verified under the calibrated harness.

\textbf{Listener-dependent modelling hurts.} Training on all 13{,}687
listener-wise ratings with learned listener embeddings and a mean-listener at
inference --- the strategy popularised by UTMOS \cite{utmos} --- reduced grouped
CV SRCC to $0.42$ / $0.44$, well below the $0.50+$ obtained from averaged labels.
With only $4.9$ ratings per pair we believe per-listener training reintroduces
the rater noise that averaging removes. This contradicts prior listener-modelling
claims and, we suspect, marks a dataset-size threshold below which the technique
is counterproductive.

\textbf{Counterfactual pairs do not decorrelate the sub-scores.} Because every
pair is a same-speaker attempt, the data never exhibits ``different speaker,
same accent''. We manufactured such pairs from natural recordings with
approximate labels, but they yielded no improvement: the metric evaluates
same-speaker pairs only, so it cannot reward decorrelation.

\textbf{Additional self-supervised backbones dilute.} Stacking HuBERT, XLS-R and
WavLM-Base alongside WavLM-Large degraded performance relative to WavLM-Large
alone.

\textbf{Phonetic posteriorgram features are redundant.} PPG features improved a
weaker configuration by $+0.013$ accent SRCC but changed the full pipeline by
only $+0.001$, suggesting that WavLM mid-layers together with CommonAccent
already capture the available pronunciation signal.

\section{Conclusion}

We presented a frozen-backbone system for pairwise speaker and accent similarity
prediction whose distinguishing components are retrieval in pair-difference space
and label-free transductive system-mean smoothing. The latter gave the largest
single gain on the official evaluation set ($+0.037$ / $+0.016$ SRCC) and lifted
our system-level correlations to third place among participating teams.

Our main limitation is calibration: because we optimised rank correlation and
rescaled outputs linearly, mean squared error was our weakest metric. Making the
predictor accurate in absolute terms without sacrificing ranking quality is the
natural next step, alongside pretraining the speaker branch on large public
voice-similarity corpora such as VoxSim \cite{voxsim}, which we identify as the
most credible remaining lever for the utterance-level gap. Code is available at
\url{https://github.com/Ranjit246/voicemos-challenge-2026-exp}.

\begin{thebibliography}{9}
\small

\bibitem{codecmos}
W.-C. Huang \emph{et al.},
``The VoiceMOS Challenge 2026 and the CodecMOS-Accent corpus,''
\emph{arXiv preprint arXiv:2603.14328}, 2026.

\bibitem{utmos}
T. Saeki \emph{et al.},
``UTMOS: UTokyo-SaruLab system for VoiceMOS Challenge 2022,''
in \emph{Proc. Interspeech}, 2022.

\bibitem{ramp}
H. Wang \emph{et al.},
``RAMP: Retrieval-augmented MOS prediction via confidence-based dynamic
weighting,'' in \emph{Proc. Interspeech}, 2023.

\bibitem{wavlm}
S. Chen \emph{et al.},
``WavLM: Large-scale self-supervised pre-training for full stack speech
processing,'' \emph{IEEE J. Sel. Topics Signal Process.}, vol.~16, no.~6, 2022.

\bibitem{commonaccent}
J. Zuluaga-Gomez \emph{et al.},
``CommonAccent: Exploring large acoustic pretrained models for accent
classification based on Common Voice,'' in \emph{Proc. Interspeech}, 2023.

\bibitem{voxsim}
J. Ahn \emph{et al.},
``VoxSim: A perceptual voice similarity dataset,''
in \emph{Proc. Interspeech}, 2024.

\bibitem{vmc23}
E. Cooper \emph{et al.},
``The VoiceMOS Challenge 2023: Zero-shot subjective speech quality prediction
for multiple domains,'' in \emph{Proc. IEEE ASRU}, 2023.

\bibitem{mosbench}
W.-C. Huang \emph{et al.},
``MOS-Bench: Benchmarking generalization ability of subjective speech quality
assessment models,'' \emph{arXiv preprint arXiv:2411.03715}, 2024.

\end{thebibliography}

\end{document}
